\chapter{数据结构}

\index{数据结构}

\key{数据结构}是一种在计算机内存中存储数据的方式。
为问题选择合适的数据结构非常重要，
因为每种数据结构都有其自身的
优势和劣势。
关键问题是：在所选的数据结构中，
哪些操作是高效的？

本章介绍C++标准库中最重要的数据结构。
尽可能使用标准库是一个好习惯，
因为它可以节省大量时间。
在本书的后半部分，我们将学习标准库中
没有提供的更高级的数据结构。

\section{动态数组}

\index{动态数组}
\index{vector}

\key{动态数组}是一种大小可以在程序执行过程中
改变的数据结构。
C++中最常用的动态数组是
\texttt{vector}结构，
它的使用方式几乎和普通数组一样。

以下代码创建一个空vector并向其中添加三个元素：

\begin{lstlisting}
vector<int> v;
v.push_back(3); // [3]
v.push_back(2); // [3,2]
v.push_back(5); // [3,2,5]
\end{lstlisting}

此后，可以像普通数组一样访问元素：

\begin{lstlisting}
cout << v[0] << "\n"; // 3
cout << v[1] << "\n"; // 2
cout << v[2] << "\n"; // 5
\end{lstlisting}

函数\texttt{size}返回vector中元素的数量。
以下代码遍历vector并打印其中的所有元素：

\begin{lstlisting}
for (int i = 0; i < v.size(); i++) {
    cout << v[i] << "\n";
}
\end{lstlisting}

\begin{samepage}
遍历vector的一种更简洁的方式如下：

\begin{lstlisting}
for (auto x : v) {
    cout << x << "\n";
}
\end{lstlisting}
\end{samepage}

函数\texttt{back}返回vector中的最后一个元素，
函数\texttt{pop\_back}移除最后一个元素：

\begin{lstlisting}
vector<int> v;
v.push_back(5);
v.push_back(2);
cout << v.back() << "\n"; // 2
v.pop_back();
cout << v.back() << "\n"; // 5
\end{lstlisting}

以下代码创建一个包含五个元素的vector：

\begin{lstlisting}
vector<int> v = {2,4,2,5,1};
\end{lstlisting}

创建vector的另一种方式是指定元素的数量
和每个元素的初始值：

\begin{lstlisting}
// 大小为10，初始值为0
vector<int> v(10);
\end{lstlisting}
\begin{lstlisting}
// 大小为10，初始值为5
vector<int> v(10, 5);
\end{lstlisting}

vector的内部实现使用的是普通数组。
如果vector的大小增加导致数组变得太小，
就会分配一个新的数组，并将所有元素
移动到新数组中。
然而，这种情况并不经常发生，
\texttt{push\_back}的平均时间复杂度为$O(1)$。

\index{string}

\texttt{string}结构
也是一种动态数组，其使用方式几乎与vector相同。
此外，字符串有一些其他数据结构所没有的特殊语法。
字符串可以使用\texttt{+}符号进行拼接。
函数$\texttt{substr}(k,x)$返回从位置$k$开始、
长度为$x$的子串，
函数$\texttt{find}(\texttt{t})$查找子串\texttt{t}第一次出现的位置。

以下代码展示了一些字符串操作：

\begin{lstlisting}
string a = "hatti";
string b = a+a;
cout << b << "\n"; // hattihatti
b[5] = 'v';
cout << b << "\n"; // hattivatti
string c = b.substr(3,4);
cout << c << "\n"; // tiva
\end{lstlisting}

\section{集合结构}

\index{set}

\key{集合}是一种维护一组元素的数据结构。
集合的基本操作包括元素
插入、查找和删除。

C++标准库包含两种集合实现：
\texttt{set}结构基于平衡二叉树，
其操作的时间复杂度为$O(\log n)$。
\texttt{unordered\_set}结构使用哈希，
其操作的平均时间复杂度为$O(1)$。

选择哪种集合实现通常是个人偏好的问题。
\texttt{set}结构的优势在于它能维护元素的顺序，
并提供\texttt{unordered\_set}中不可用的函数。
另一方面，\texttt{unordered\_set}
可能更高效。

以下代码创建一个包含整数的集合，
并展示了一些操作。
函数\texttt{insert}向集合中添加元素，
函数\texttt{count}返回元素在集合中出现的次数，
函数\texttt{erase}从集合中移除元素。

\begin{lstlisting}
set<int> s;
s.insert(3);
s.insert(2);
s.insert(5);
cout << s.count(3) << "\n"; // 1
cout << s.count(4) << "\n"; // 0
s.erase(3);
s.insert(4);
cout << s.count(3) << "\n"; // 0
cout << s.count(4) << "\n"; // 1
\end{lstlisting}

集合的使用方式大体上与vector类似，
但不能使用\texttt{[]}语法来访问元素。
以下代码创建一个集合，
打印其中元素的数量，然后遍历所有元素：
\begin{lstlisting}
set<int> s = {2,5,6,8};
cout << s.size() << "\n"; // 4
for (auto x : s) {
    cout << x << "\n";
}
\end{lstlisting}

集合的一个重要性质是
所有元素都是\emph{互异的}。
因此，函数\texttt{count}总是返回
0（元素不在集合中）
或1（元素在集合中），
而函数\texttt{insert}在元素已存在时
不会再次添加。
以下代码说明了这一点：

\begin{lstlisting}
set<int> s;
s.insert(5);
s.insert(5);
s.insert(5);
cout << s.count(5) << "\n"; // 1
\end{lstlisting}

C++还包含\texttt{multiset}和\texttt{unordered\_multiset}结构，
它们的使用方式与\texttt{set}
和\texttt{unordered\_set}类似，
但它们可以包含同一个元素的多个实例。
例如，以下代码将数字5的三个实例
都添加到一个multiset中：

\begin{lstlisting}
multiset<int> s;
s.insert(5);
s.insert(5);
s.insert(5);
cout << s.count(5) << "\n"; // 3
\end{lstlisting}
函数\texttt{erase}会从multiset中
移除一个元素的所有实例：
\begin{lstlisting}
s.erase(5);
cout << s.count(5) << "\n"; // 0
\end{lstlisting}
通常，只需要移除一个实例，
可以按如下方式实现：
\begin{lstlisting}
s.erase(s.find(5));
cout << s.count(5) << "\n"; // 2
\end{lstlisting}

\section{映射结构}

\index{map}

\key{映射}是一种广义数组，
由键值对组成。
普通数组的键总是连续的整数$0,1,\ldots,n-1$，
其中$n$是数组的大小，
而映射的键可以是任何数据类型，
且不需要是连续的值。

C++标准库包含两种映射实现，
与集合的实现相对应：
\texttt{map}结构基于平衡二叉树，
访问元素的时间复杂度为$O(\log n)$，
而\texttt{unordered\_map}结构使用哈希，
访问元素的平均时间复杂度为$O(1)$。

以下代码创建一个键为字符串、值为整数的映射：

\begin{lstlisting}
map<string,int> m;
m["monkey"] = 4;
m["banana"] = 3;
m["harpsichord"] = 9;
cout << m["banana"] << "\n"; // 3
\end{lstlisting}

如果请求一个键的值，
但映射中不包含该键，
则该键会自动添加到映射中并赋予默认值。
例如，在以下代码中，
键''aybabtu''的值0被添加到映射中。

\begin{lstlisting}
map<string,int> m;
cout << m["aybabtu"] << "\n"; // 0
\end{lstlisting}
函数\texttt{count}检查映射中
是否存在某个键：
\begin{lstlisting}
if (m.count("aybabtu")) {
    // 键存在
}
\end{lstlisting}
以下代码打印映射中所有的键和值：
\begin{lstlisting}
for (auto x : m) {
    cout << x.first << " " << x.second << "\n";
}
\end{lstlisting}

\section{迭代器与范围}

\index{迭代器}

C++标准库中的许多函数都使用迭代器来操作。
\key{迭代器}是一个指向
数据结构中某个元素的变量。

常用的迭代器\texttt{begin}
和\texttt{end}定义了一个范围，
该范围包含数据结构中的所有元素。
迭代器\texttt{begin}指向数据结构中的第一个元素，
而迭代器\texttt{end}指向最后一个元素\emph{之后}的位置。
情况如下图所示：

\begin{center}
\begin{tabular}{llllllllll}
\{ & 3, & 4, & 6, & 8, & 12, & 13, & 14, & 17 & \} \\
& $\uparrow$ & & & & & & & & $\uparrow$ \\
& \multicolumn{3}{l}{\texttt{s.begin()}} & & & & & & \texttt{s.end()} \\
\end{tabular}
\end{center}

注意迭代器中的不对称性：
\texttt{s.begin()}指向数据结构中的一个元素，
而\texttt{s.end()}指向数据结构之外。
因此，迭代器定义的范围是\emph{半开}的。

\subsubsection{范围操作}

迭代器用于C++标准库中那些需要接收数据结构中
一个元素范围的函数。
通常，我们希望处理数据结构中的所有元素，
因此需要为函数提供\texttt{begin}和\texttt{end}迭代器。

例如，以下代码使用函数\texttt{sort}对一个vector排序，
然后使用函数\texttt{reverse}反转元素的顺序，
最后使用函数\texttt{random\_shuffle}随机打乱元素的顺序。

\index{sort@\texttt{sort}}
\index{reverse@\texttt{reverse}}
\index{random\_shuffle@\texttt{random\_shuffle}}

\begin{lstlisting}
sort(v.begin(), v.end());
reverse(v.begin(), v.end());
random_shuffle(v.begin(), v.end());
\end{lstlisting}

这些函数也可以用于普通数组。
在这种情况下，需要为数组提供指针
而不是迭代器：

\newpage
\begin{lstlisting}
sort(a, a+n);
reverse(a, a+n);
random_shuffle(a, a+n);
\end{lstlisting}

\subsubsection{集合迭代器}

迭代器常用于访问集合中的元素。
以下代码创建一个迭代器\texttt{it}，
指向集合中最小的元素：
\begin{lstlisting}
set<int>::iterator it = s.begin();
\end{lstlisting}
一种更简洁的写法如下：
\begin{lstlisting}
auto it = s.begin();
\end{lstlisting}
可以使用\texttt{*}符号访问迭代器所指向的元素。
例如，以下代码打印集合中的第一个元素：

\begin{lstlisting}
auto it = s.begin();
cout << *it << "\n";
\end{lstlisting}

可以使用运算符\texttt{++}（向前）和\texttt{--}（向后）
来移动迭代器，
即迭代器会移动到集合中的下一个或上一个元素。

以下代码按递增顺序打印所有元素：
\begin{lstlisting}
for (auto it = s.begin(); it != s.end(); it++) {
    cout << *it << "\n";
}
\end{lstlisting}
以下代码打印集合中最大的元素：
\begin{lstlisting}
auto it = s.end(); it--;
cout << *it << "\n";
\end{lstlisting}

函数$\texttt{find}(x)$返回一个
指向值为$x$的元素的迭代器。
但如果集合中不包含$x$，
迭代器将为\texttt{end}。

\begin{lstlisting}
auto it = s.find(x);
if (it == s.end()) {
    // 未找到x
}
\end{lstlisting}

函数$\texttt{lower\_bound}(x)$返回
一个迭代器，指向集合中值\emph{至少为}$x$的最小元素，
函数$\texttt{upper\_bound}(x)$
返回一个迭代器，指向集合中值\emph{大于}$x$的最小元素。
在两个函数中，如果这样的元素不存在，
返回值均为\texttt{end}。
这两个函数不被\texttt{unordered\_set}结构支持，
因为该结构不维护元素的顺序。

\begin{samepage}
例如，以下代码查找最接近$x$的元素：

\begin{lstlisting}
auto it = s.lower_bound(x);
if (it == s.begin()) {
    cout << *it << "\n";
} else if (it == s.end()) {
    it--;
    cout << *it << "\n";
} else {
    int a = *it; it--;
    int b = *it;
    if (x-b < a-x) cout << b << "\n";
    else cout << a << "\n";
}
\end{lstlisting}

该代码假设集合不为空，
并使用迭代器\texttt{it}遍历所有可能的情况。
首先，迭代器指向值至少为$x$的最小元素。
如果\texttt{it}等于\texttt{begin}，
则对应的元素最接近$x$。
如果\texttt{it}等于\texttt{end}，
则集合中最大的元素最接近$x$。
如果以上情况都不成立，
最接近$x$的元素要么是\texttt{it}对应的元素，
要么是前一个元素。
\end{samepage}

\section{其他结构}

\subsubsection{Bitset}

\index{bitset}

\key{bitset（位集）}是一种每个值
只能是0或1的数组。
例如，以下代码创建一个包含10个元素的bitset：
\begin{lstlisting}
bitset<10> s;
s[1] = 1;
s[3] = 1;
s[4] = 1;
s[7] = 1;
cout << s[4] << "\n"; // 1
cout << s[5] << "\n"; // 0
\end{lstlisting}

使用bitset的优势在于
它们比普通数组占用更少的内存，
因为bitset中的每个元素只占用一位内存。
例如，
如果将$n$个位存储在\texttt{int}数组中，
将使用$32n$位内存，
而相应的bitset只需$n$位内存。
此外，bitset的值可以使用位运算高效操作，
这使得用位集优化算法成为可能。

以下代码展示了创建上述bitset的另一种方式：
\begin{lstlisting}
bitset<10> s(string("0010011010")); // 从右到左
cout << s[4] << "\n"; // 1
cout << s[5] << "\n"; // 0
\end{lstlisting}

函数\texttt{count}返回
bitset中1的个数：

\begin{lstlisting}
bitset<10> s(string("0010011010"));
cout << s.count() << "\n"; // 4
\end{lstlisting}

以下代码展示了位运算的使用示例：
\begin{lstlisting}
bitset<10> a(string("0010110110"));
bitset<10> b(string("1011011000"));
cout << (a&b) << "\n"; // 0010010000
cout << (a|b) << "\n"; // 1011111110
cout << (a^b) << "\n"; // 1001101110
\end{lstlisting}

\subsubsection{Deque}

\index{deque}

\key{deque（双端队列）}是一种动态数组，
其大小可以在数组的两端高效地改变。
与vector一样，deque提供了函数
\texttt{push\_back}和\texttt{pop\_back}，
但它还包含了vector中不可用的函数
\texttt{push\_front}和\texttt{pop\_front}。

deque的使用方式如下：
\begin{lstlisting}
deque<int> d;
d.push_back(5); // [5]
d.push_back(2); // [5,2]
d.push_front(3); // [3,5,2]
d.pop_back(); // [3,5]
d.pop_front(); // [5]
\end{lstlisting}

deque的内部实现比vector更复杂，
因此deque比vector更慢。
不过，在两端添加和移除元素的平均时间复杂度均为$O(1)$。

\subsubsection{Stack}

\index{stack}

\key{stack（栈）}
是一种提供两种$O(1)$时间操作的数据结构：
在顶部添加元素，
以及从顶部移除元素。
只能访问栈顶的元素。

以下代码展示了如何使用栈：
\begin{lstlisting}
stack<int> s;
s.push(3);
s.push(2);
s.push(5);
cout << s.top(); // 5
s.pop();
cout << s.top(); // 2
\end{lstlisting}
\subsubsection{Queue}

\index{queue}

\key{queue（队列）}也提供两种$O(1)$时间操作：
在队列末尾添加元素，
以及移除队列中的第一个元素。
只能访问队列的第一个和最后一个元素。

以下代码展示了如何使用队列：
\begin{lstlisting}
queue<int> q;
q.push(3);
q.push(2);
q.push(5);
cout << q.front(); // 3
q.pop();
cout << q.front(); // 2
\end{lstlisting}

\subsubsection{优先队列}

\index{优先队列}
\index{堆}

\key{优先队列}（priority queue）
维护一组元素。
支持的操作包括插入，以及根据队列类型
检索和移除最小或最大元素。
插入和移除的时间复杂度为$O(\log n)$，
检索的时间复杂度为$O(1)$。

虽然有序集合也能高效支持优先队列的所有操作，
但使用优先队列的优势在于它的常数因子更小。
优先队列通常使用堆结构实现，
这比有序集合中使用的平衡二叉树简单得多。

\begin{samepage}
默认情况下，C++优先队列中的元素
按递减顺序排列，
可以查找和移除队列中最大的元素。
以下代码说明了这一点：

\begin{lstlisting}
priority_queue<int> q;
q.push(3);
q.push(5);
q.push(7);
q.push(2);
cout << q.top() << "\n"; // 7
q.pop();
cout << q.top() << "\n"; // 5
q.pop();
q.push(6);
cout << q.top() << "\n"; // 6
q.pop();
\end{lstlisting}
\end{samepage}

如果我们想创建一个支持查找和移除
最小元素的优先队列，
可以按如下方式实现：

\begin{lstlisting}
priority_queue<int,vector<int>,greater<int>> q;
\end{lstlisting}

\subsubsection{基于策略的数据结构}

\texttt{g++}编译器还支持
一些不属于C++标准库的数据结构。
这些结构被称为\emph{基于策略的}数据结构。
要使用这些结构，必须在代码中添加以下行：
\begin{lstlisting}
#include <ext/pb_ds/assoc_container.hpp>
using namespace __gnu_pbds; 
\end{lstlisting}
此后，我们可以定义一个数据结构\texttt{indexed\_set}，
它类似于\texttt{set}，但可以像数组一样进行索引。
对于\texttt{int}值的定义如下：
\begin{lstlisting}
typedef tree<int,null_type,less<int>,rb_tree_tag,
             tree_order_statistics_node_update> indexed_set; 
\end{lstlisting}
现在我们可以按如下方式创建一个集合：
\begin{lstlisting}
indexed_set s;
s.insert(2);
s.insert(3);
s.insert(7);
s.insert(9);
\end{lstlisting}
这个集合的特殊之处在于，我们可以访问
元素在有序数组中对应的索引。
函数$\texttt{find\_by\_order}$返回
指向给定位置元素的迭代器：
\begin{lstlisting}
auto x = s.find_by_order(2);
cout << *x << "\n"; // 7
\end{lstlisting}
函数$\texttt{order\_of\_key}$
返回给定元素的位置：
\begin{lstlisting}
cout << s.order_of_key(7) << "\n"; // 2
\end{lstlisting}
如果元素不在集合中，
我们得到的是该元素在集合中
将会处于的位置：
\begin{lstlisting}
cout << s.order_of_key(6) << "\n"; // 2
cout << s.order_of_key(8) << "\n"; // 3
\end{lstlisting}
这两个函数都以对数时间运行。

\section{与排序的比较}

通常可以用数据结构或排序来解决同一个问题。
有时这两种方法的实际效率
存在显著差异，
而这种差异可能隐藏在它们的时间复杂度中。

考虑一个问题：给定两个列表$A$和$B$，
它们都包含$n$个元素。
我们的任务是计算同时属于两个列表的元素数量。
例如，对于列表
\[A = [5,2,8,9] \hspace{10pt} \textrm{和} \hspace{10pt} B = [3,2,9,5],\]
答案是3，因为数字2、5和9
同时属于两个列表。

该问题的一个直接解决方案是
在$O(n^2)$时间内遍历所有元素对，
但接下来我们将关注更高效的算法。

\subsubsection{算法1}

我们用$A$中出现的元素构造一个集合，
然后遍历$B$的元素，
检查每个元素是否也属于$A$。
这种方法很高效，因为$A$的元素存储在集合中。
使用\texttt{set}结构，
算法的时间复杂度为$O(n \log n)$。

\subsubsection{算法2}

没有必要维护一个有序集合，
因此我们可以用\texttt{unordered\_set}结构
代替\texttt{set}结构。
这是一种使算法更高效的简单方法，
因为我们只需要更改底层数据结构。
新算法的时间复杂度为$O(n)$。

\subsubsection{算法3}

除了数据结构，我们还可以使用排序。
首先，对两个列表$A$和$B$分别排序。
然后，同时遍历两个列表，
找出共同的元素。
排序的时间复杂度为$O(n \log n)$，
算法的其余部分在$O(n)$时间内运行，
因此总时间复杂度为$O(n \log n)$。

\subsubsection{效率比较}

下表展示了当$n$变化且列表元素为$1 \ldots 10^9$之间的随机整数时，
上述算法的效率：

\begin{center}
\begin{tabular}{rrrr}
$n$ & 算法1 & 算法2 & 算法3 \\
\hline
$10^6$ & $1.5$ s & $0.3$ s & $0.2$ s \\
$2 \cdot 10^6$ & $3.7$ s & $0.8$ s & $0.3$ s \\
$3 \cdot 10^6$ & $5.7$ s & $1.3$ s & $0.5$ s \\
$4 \cdot 10^6$ & $7.7$ s & $1.7$ s & $0.7$ s \\
$5 \cdot 10^6$ & $10.0$ s & $2.3$ s & $0.9$ s \\
\end{tabular}
\end{center}

算法1和算法2的唯一区别在于
它们使用了不同的集合结构。
在这个问题中，这个选择对运行时间有重要影响，
因为算法2比算法1快4–5倍。

然而，最有效的算法是使用排序的算法3。
它所需的时间仅为算法2的一半。
有趣的是，算法1和算法3的时间复杂度
都是$O(n \log n)$，
尽管如此，算法3还是快了十倍。
这可以解释为：排序是一个简单的操作，
只在算法3的开头执行一次，
而算法的其余部分以线性时间运行。
另一方面，
算法1在整个算法过程中都要维护
一个复杂的平衡二叉树。
